\documentclass[12pt,a4paper]{report}

\usepackage[utf8]{inputenc} % for diacritics support
\usepackage[english]{babel} % settings for the English language
\renewcommand\familydefault{\sfdefault} % sans serif

\usepackage[margin=2.54cm]{geometry}	% page dimensions and margins
\usepackage{graphicx} % support the \includegraphics command and options

% formatting sections and subsections
\usepackage{textcase}
\usepackage[titletoc, title]{appendix}
\usepackage{titlesec}
\titleformat{\chapter}{\large\bfseries\MakeUppercase}{\thechapter}{2ex}{}[\vspace*{-1.5cm}]
\titleformat*{\section}{\large\bfseries}
\titleformat*{\subsection}{\large\bfseries}
\titleformat*{\subsubsection}{\large\bfseries}

\usepackage{chngcntr}
\counterwithout{figure}{chapter} % no chapter number in figure labels
\counterwithout{table}{chapter} % no chapter number in table labels
\counterwithout{equation}{chapter} % no chapter number in equation labels

\usepackage{booktabs} % for much better looking tables
\usepackage{url} % Useful for inserting web links nicely
\usepackage[bookmarks,unicode,hidelinks]{hyperref}

\usepackage{array} % for better arrays (eg matrices) in maths
\usepackage{paralist} % very flexible & customisable lists (eg. enumerate/itemize, etc.)
\usepackage{verbatim} % adds environment for commenting out blocks of text & for better verbatim
\usepackage{subfig} % make it possible to include more than one captioned figure/table in a single float
\usepackage{enumitem}
\setlist{noitemsep}

%%% HEADERS & FOOTERS
\usepackage{fancyhdr}
\pagestyle{empty}
\renewcommand{\headrulewidth}{0pt}
\renewcommand{\footrulewidth}{0pt}
\lhead{}\chead{}\rhead{}
\lfoot{}\cfoot{\thepage}\rfoot{}



\newcommand{\HeaderLineSpace}{-0.25cm}
\newcommand{\UniTextEN}{UNIVERSITY POLITEHNICA OF BUCHAREST \\[\HeaderLineSpace]
FACULTY OF AUTOMATIC CONTROL AND COMPUTERS \\[\HeaderLineSpace]
COMPUTER SCIENCE AND ENGINEERING DEPARTMENT\\}
\newcommand{\DiplomaEN}{DIPLOMA PROJECT}
\newcommand{\AdvisorEN}{Thesis advisor:}
\newcommand{\BucEN}{BUCHAREST}

\newcommand{\frontPage}[6]{
\begin{titlepage}
\begin{center}
{\Large #1}  % header (university, faculty, department)
\vspace{50pt}
\begin{tabular}{p{6cm}p{4cm}}
\includegraphics[scale=0.8]{pics/upb-logo.jpg} &
	\includegraphics[scale=0.5,trim={14cm 11cm 2cm 5cm},clip=true]{pics/cs-logo.pdf}
\end{tabular}

\vspace{105pt}
{\Huge #2}\\                           % diploma project text
\vspace{40pt}
{\Large #3}\\ \vspace{0pt}  % project title
{\Large #4}\\                          % project subtitle
\vspace{40pt}
{\LARGE \Name}\\                   % student name
\end{center}
\vspace{60pt}
\begin{tabular*}{\textwidth}{@{\extracolsep{\fill}}p{6cm}r}
&{\large\textbf{#5}}\vspace{10pt}\\      % scientific advisor
&{\large \Advisor}                                    % advisor name
\end{tabular*}
\vspace{20pt}
\begin{center}
{\large\textbf{#6}}\\                                % bucharest
\vspace{0pt}
{\normalsize \Year}
\end{center}
\end{titlepage}
}

\newcommand{\frontPageEN}{\frontPage{\UniTextEN}{\DiplomaEN}{\ProjectTitleEN}{\ProjectSubtitleEN}{\AdvisorEN}{\BucEN}}

\linespread{1.15}
\setlength\parindent{0pt}
\setlength\parskip{.28cm}

%% Abstract macro
\newcommand{\AbstractPage}{
\begin{titlepage}
\textbf{\large ABSTRACT}\par
\AbstractEN \vfill
\end{titlepage}
}

%% Thank you macro
\newcommand{\ThanksPage}{
\begin{titlepage}
{\noindent \large\textbf{ACKNOWLEDGEMENTS}}\\
\Thanks
\end{titlepage}
}



%%%%%%%%%%%%%%%%%%%%%%%%%%%%%%%%%%%%%%%%%%%%%%%%%%
%%
%%          End of template definitions
%%
%%%%%%%%%%%%%%%%%%%%%%%%%%%%%%%%%%%%%%%%%%%%%%%%%%


%%% You can remove these lines from your thesis; they are used only by the template.
\newcommand{\worktype}[1]{[\textit{#1}] }
\newcommand{\dezvoltare}{\worktype{Product development}}
\newcommand{\cercetare}{\worktype{Research}}
\newcommand{\ambele}{\worktype{Both}}
%%%


%%
%%   The fields below must be modified by the author. Change only the content, not the name of each definition.
%%
\newcommand{\ProjectTitleEN}{Diploma Project Title  (eg: Diploma project template)}
\newcommand{\ProjectSubtitleEN}{Subtitle (eg: 2018 version)}
\newcommand{\Name}{Ioana Popescu}
\newcommand{\Advisor}{Prof. dr. ing. Andrei Ionescu}
\newcommand{\Year}{2018}

% Document settings
\title{Diploma Project}
\author{\Name}
\date{\Year}

%%
%%   Abstract fields
%%
\newcommand{\AbstractEN}{The abstract has an introductory role and should engulf both a brief description of the issue at hand, as well as an overview of the obtained results and conclusions. The abstract should be formulated such that even somebody that is unfamiliar with the projects’ domain can grasp the objectives of the thesis while, at the same time, retaining a specificity level offering a bird’s eye view of the project.
The recommended size for this section is limited to 200 words.}

%%
%%   Acknowledgements page fields
%%
\newcommand{\Thanks}{(optional) Here you can add a special acknowledgements section. }

\begin{document}

\frontPageEN

\begingroup
\linespread{1}
\tableofcontents
\endgroup

\AbstractPage

% can be commented out or deleted
\ThanksPage


% The thesis text begins here



\chapter{Introduction}\pagestyle{fancy}
% * <marios.choudary@gmail.com> 2018-02-28T11:38:18.106Z:
%
% > INTRODUCTION
% I removed the font references here so that we are no longer dependent on Calibri. Personally, I am not even sure this recommendation helps much, and the default LaTeX font (Computer Modern) seems fine to me. If you agree, please delete the commented-out lines below, as well as the ones referring to the Calibri font in the rest of the document.
%
% ^.
Recommended formatting parameters for the thesis:
\begin{itemize}
 %\item Recommended font: Calibri; Font size: 12;
 \item Font size: 12;
 \item Line spacing: 1.5; Spacing after paragraph: 8pt;
 \item Style: Justified;
 \item Page size: A4; Margins: 2.54cm/ 2.54cm/ 2.54cm/ 2.54cm;
 %\item Heading1: Calibri, 14, bold, all caps;
 %\item Heading2: Calibri, 14, bold;
 %\item Heading3: Calibri, 12.
 %\item Font for formulas: Cambria Math, 12.
 \item Heading1: 14, bold, all caps;
 \item Heading2: 14, bold;
 \item Heading3: 12.
 \item Font size for formulas: 12.
\end{itemize}
In the introduction, it is necessary to address the following points, which essentially familiarize the reader (the committee, fellow students, or domain experts) with the topic of the project, the proposed solution, and the contents/structure of the thesis. Although the introduction may also contain some more general elements, it is recommended to maintain a technical language, specific to the audience that will read the thesis.

In the following chapters, you will find a series of notations of the form \dezvoltare, \cercetare. This type of formatting is used exclusively in this template to mark advice and requirements specific to diploma theses of different kinds. When preparing your own document, you will not use these markers.
The elements you must address in the introduction are described in the subsections below.
\section{Context}
A short introduction to the project, the motivation, and an explanation of why the project's domain is relevant.
\section{Problem}
What is the problem that the project will solve.
\section{Objectives}
What are the objectives of the project/solution/approach/idea; what improvements or developments result from solving the project.
\section{Proposed Solution}
A brief description of the implemented solution; what approach is proposed (not the details of tools and technologies, but the approach and idea proposed by the author).
\section{Obtained Results}
A brief description of the obtained results, and possibly why they are important compared to other solutions or studies.
\section{Thesis Structure}
A paragraph in which each of the following sections is presented in 1-2 sentences, emphasizing the most significant elements of each section.



\chapter{Requirements Analysis / Motivation}
\dezvoltare This chapter will analyze the product requirements from the perspective of potential customers and the anticipated usage scenarios, resulting in a list of features.

\cercetare This chapter will introduce the motivation for carrying out the proposed project.

If the bachelor's project is part of a larger project (for example, a complex project worked on by 2 students (e.g.: 1 student on the application's front-end, 1 student on the application's back-end)), this chapter will briefly explain the overall project and which part of the project is addressed by the proposed thesis.

Criteria for the \textit{Un\textit{satisfactory}} grade:
\begin{itemize}
	\item \dezvoltare The requirements are imagined by the student based on a market analysis;
	\item \cercetare No valid motivation is provided.
\end{itemize}

Criteria for the \textit{Satisfactory} grade:
\begin{itemize}
	\item \dezvoltare There is an interview, a client, and the requirements analysis is developed based on the interview;
	\item \cercetare The motivation is only personal.
\end{itemize}


Criteria for the \textit{Good} grade:
\begin{itemize}
	\item	 \dezvoltare An iterative process based on interviews with several clients, MVP development, re-evaluation of requirements;
	\item	 \cercetare The motivation is linked to an explicit scientific / technical need.
\end{itemize}


\chapter{Market Study / Existing Methods}
\dezvoltare What similar solutions exist on the market? What are their limitations / for which use cases or for which type of customers do the products available on the market fail to meet the requirements? What are the indicators by which these products are evaluated by potential customers, and where are the gaps / what is the opportunity created by these gaps?

\cercetare Existing methods (or ``State of the Art'') usually refer to the current level of development: what is the current state of the field, where we stand, what the context is. What are the current solutions present in the specialized literature and what are their limitations? What exploration directions are recommended in the specialized literature? The specialized literature refers to recent scientific articles, published in journals with a high impact factor, or in the proceedings of top conferences, or in books.

\ambele At the end of this chapter, the aim is to describe the technologies used in the thesis, together with alternatives and convincing qualitative and quantitative arguments.

Criteria for the \textit{Un\textit{satisfactory}} grade:
\begin{itemize}
	\item \dezvoltare A few products on the market are analyzed superficially;
	\item \cercetare the literature review is limited to research groups from Romania;
	\item \ambele The technologies used in the thesis are described.
\end{itemize}

Criteria for the \textit{Satisfactory} grade:
\begin{itemize}
	\item \dezvoltare There is an interview, a client, and the requirements analysis is developed based on the interview.
	\item \cercetare a review of the international specialized literature, without precisely positioning the thesis within the landscape of the studied field;
	\item \ambele A few alternative technologies are described for each of the technologies used in the thesis. There is a justification regarding the choice.
\end{itemize}

Criteria for the \textit{Good} grade:
\begin{itemize}
	\item \dezvoltare An iterative process based on interviews with several clients, MVP development, re-evaluation of requirements;
	\item \cercetare a review of the international specialized literature, precisely positioning the thesis within the current landscape of the studied field;
	\item \ambele Alternative technologies are described. They are analyzed quantitatively and qualitatively, using benchmarks and tests performed by the student. The analysis is summarized through tables and graphs.
\end{itemize}

\section{Figure Formatting Guidelines}

The figures used in the document will be centered and numbered (for example Figure~\ref{fig:pic1}).
Any figure that is not created by the author of the thesis must mandatorily be cited either at the end (for example Figure ~\ref{fig:pic2} is taken from the document \cite{article}), or at least in a footnote (see Figure~\ref{fig:pic2}). Any figure exceeding 50\% of a page in size will be moved to the appendices. All figures in the thesis will be referenced in the text. Example: Figure~\ref{fig:pic1} shows a schematic diagram of an inverting amplifier with an op-amp.

\begin{figure}[th]
\centering
\includegraphics{pics/Pic1.png}
  \caption{Inverting amplifier}
  \label{fig:pic1}
\end{figure}

\newpage

\begin{figure}[th]
\centering
\includegraphics{pics/Pic2.png}
  \caption[Instrumentation amplifier with 3 op-amps]{Instrumentation amplifier with 3 op-amps\protect\footnotemark}
  \label{fig:pic2}
\end{figure}
\footnotetext{© http://www.ece.tamu.edu/sspalermo/ecen3205/Secton\%201III.pdf}

\chapter{Proposed Solution}
This chapter contains an overview of the solution that solves the problem, by presenting its structure / architecture. Depending on the type of thesis, this chapter may contain diagrams (class, deployment, workflow, entity-relationship), correctness proofs for the algorithms proposed by the author, theoretical approaches (mathematical modeling), the hardware structure, and the application architecture.


Criteria for the \textit{Un\textit{satisfactory}} grade:
\begin{itemize}
	\item	Description in natural language.
\end{itemize}

Criteria for the \textit{Satisfactory} grade:
\begin{itemize}
	\item	Description + database, workflow, class, and algorithm diagrams.
\end{itemize}

Criteria for the \textit{Good} grade:
\begin{itemize}
	\item 	Description + database, workflow, class, and algorithm diagrams + a description of the process through which the architecture/structure of the solution was developed.
\end{itemize}

\section{Formula Formatting Guidelines}
The mathematical formulas used in the document will be centered on the page and numbered.

\begin{equation}
(x+a)^n = \sum_{k=0}^{n}\left(\begin{array}{c}n\\k\\\end{array}\right)x^ka^{n-k}
\end{equation}

\begin{equation}
f(x) = a_0 + \sum_{n=1}^{\infty}\left(a_n \cos\frac{n\pi x}{L} + b_n\sin\frac{n\pi x}{L}\right)
\end{equation}



\chapter{Implementation Details}
In addition to the previous chapter, this one contains specific elements of solving the problem that involved particular technical difficulties. It may include configurations, code snippets, pseudo-code, implementations of algorithms, data analyses, and test scripts. It may also detail how the technologies introduced in Chapter 3 were used.


Criteria for the \textit{Un\textit{satisfactory}} grade:
\begin{itemize}
	\item	Diagrams and pseudo-code are briefly presented.
\end{itemize}
Criteria for the \textit{Satisfactory} grade:
\begin{itemize}
	\item	A summary description of the implementation, the presentation of irrelevant code snippets, diagrams, etc.
\end{itemize}
Criteria for the \textit{Good} grade:
\begin{itemize}
	\item	A detailed description of the algorithms/structures used; A step-by-step presentation of the development, including implementation difficulties encountered and solutions discovered; (where applicable) a proof of correctness of the algorithms used.
\end{itemize}

\section{Table Formatting Guidelines}
It is recommended to use tables of the form shown below.  Font size:  9.
Any table present in the thesis will be referenced in the text; example: see Table~\ref{tab:criterii}.

\begin{table}[th]\small\linespread{1}
\caption{Summary of criteria}
\label{tab:criterii}
\begin{tabular}{l >{\raggedright\arraybackslash}p{8cm} >{\raggedright\arraybackslash}p{4cm}}
\textbf{Grade} & \textbf{Criterion} & \textbf{Notes} \\\hline
\textbf{Unsatisfactory} & Diagrams and pseudo-code are briefly presented & \\\hline
\textbf{Satisfactory} &A summary description of the implementation, the presentation of irrelevant code snippets, diagrams, etc.& \\
\hline
\textbf{\textit{Good}} &A detailed description of the algorithms/structures used; A step-by-step presentation of the development, including implementation difficulties encountered and solutions discovered; (where applicable) a proof of correctness of the algorithms used. & May include configurations, code snippets, pseudo-code, implementations of algorithms, data analyses, and test scripts. \\
\hline
\end{tabular}
\end{table}


\chapter{Evaluation}
This chapter must, in principle, answer 2 questions and conclude with a discussion of the obtained results. The two questions that must be answered are:
\begin{enumerate}
	\item  \textbf{Does it work correctly?} (According to the specifications extracted in Chapter 2);
The evaluation of whether it works correctly is done based on the requirements identified in the previous chapters.

	\item How \textit{well} does it work / how does it compare to existing solutions? (based on clear metrics).
The evaluation of how \textit{well} it works must be based on percentages, times, quantities, numbers, \textbf{compared to the solutions presented in Chapter 3}. This may concern performance, overhead, consumed resources, scalability, etc.
\end{enumerate}

When carrying out the discussion, tables with percentages, numerical results, and graphs will be used. Usually, comparisons and comparative tests with other similar projects (if they exist) are made here, and strengths and weaknesses are extracted. The mentioned advantages are taken into account and the viability of the approach / application is demonstrated, preferably through comparison with other approaches (if this is possible). The key word in evaluation is ``metric'': you must have measurable and quantifiable notions. During the evaluation process, explain the data, tables, and graphs that you present and emphasize their relevance, in the following style: ``it is preferable ... because …''; explain to the reader not only the data but also its meaning and how it is interpreted. From this interpretation, the positioning of your project among the existing alternatives must result, as well as how it can be further improved.

Criteria for the \textit{Un\textit{satisfactory}} grade:
\begin{itemize}
	\item The application is tested but runs on the student's computer, there are no testing facilities, and it has not been validated with clients / users;
	\item No comparisons with other similar systems have been made.
\end{itemize}

Criteria for the \textit{Satisfactory} grade:
\begin{itemize}
	\item \dezvoltare  There are unit and integration tests, there is a deployment strategy, and there is minimal validation with clients / users.
	\item \cercetare The main components and the overall solution have been evaluated in terms of performance, but standard datasets are not used and there are some errors in interpreting the data.
	\item \ambele A minimal discussion of the relevance of the presented results, and a minimal comparison with other similar systems.
\end{itemize}

Criteria for the \textit{Good} grade:
\begin{itemize}
	\item \dezvoltare Unit and integration tests, deployment tools that are used and that show consistent work throughout the semester, a thesis validated with clients / users, and a product in production.
	\item \cercetare The components and the overall solution have been evaluated in terms of performance, using standard datasets and with a correct interpretation of the results.
	\item \ambele A discussion with a qualitative and quantitative presentation of the results, as well as their relevance, through a comprehensive comparison with other similar systems.
\end{itemize}

\chapter{Conclusions}
This chapter summarizes the entire project, from objectives, to implementation, and to the relevance of the obtained results. At the end of the chapter, there may be a ``Future Work'' subsection.

Criteria for the \textit{Un\textit{satisfactory}} grade:
\begin{itemize}
	\item	The conclusions are not correlated with the content of the thesis;
\end{itemize}

Criteria for the \textit{Satisfactory} grade:
\begin{itemize}
	\item	The conclusions are correlated with the content of the thesis, but they do not provide a picture of the quality and relevance of the obtained results;
\end{itemize}

Criteria for the \textit{Good} grade:
\begin{itemize}
	\item	The conclusions are correlated with the content of the thesis, and they provide a precise picture of the relevance and quality of the results obtained within the project.
\end{itemize}

\chapter*{Bibliography}\addcontentsline{toc}{chapter}{Bibliography}
% * <marios.choudary@gmail.com> 2018-02-28T12:07:48.730Z:
%
% > BIBLIOGRAPHY
% I added a paragraph with a few details about using bibliographic citations in LaTeX, about using ``\cite'', and about the possibility of including the bibliography directly in the LaTeX file.
%
% ^.

\begin{itemize}
	\item 	Do NOT use references to Wikipedia or other sources without an acknowledged author.
	\item 	For references to relevant articles accessible on the web (described by URL), the access date will also be noted in the bibliography.
	\item 	More details about citing internet references can be found at:
	\begin{itemize}
		\item	\url{http://www.writinghelp-central.com/apa-citation-internet.html}
		\item	\url{http://www.webliminal.com/search/search-web13.html}
	\end{itemize}
	\item 	Footnotes are used if you reference a less significant link a single time; if the note is cited several times, then use a bibliographic reference.
	\item 	If an image is inserted into the text and is not created by the author of the thesis, its source must be cited (as a footnote or reference - using a footnote is preferable).
	\item 	References are placed directly connected to the text (for example ``KVM [1] uses'', ``as stated by Popescu and Ionescu [12]'', etc.). It is not recommended to use formulations such as ``[1] uses'', ``as stated in [12]'', ``as described in [11]'', etc.
	\item 	Statements of the form ``are numerous'', ``have grown exponentially'', ``are among the most used'', ``are an important topic'' must be backed by citations, concrete data, and comparative analyses.
	\begin{itemize}
		\item	Especially in the introduction, ``state of the art'', ``related work'', or ``background'' chapters, you must support your statements with citations. Be self-critical and consider whether your statements need citations, even those you consider obvious.
		\item	Most of the citations will be in the introduction, ``state of the art'', ``related work'', or ``background'' chapters.
	\end{itemize}
	\item 	All bibliographic entries must be cited in the text. Do not simply add them at the end.
	\item 	Never copy or translate from information sources of any kind (online, offline, books, etc.). If you nevertheless wish to provide, by exception, a famous quotation - of at most 1 sentence - use quotation marks and obviously mention the source.
	\item 	If you reformulate ideas or create a paragraph summarizing some ideas in your own words, indicate with a citation (bibliographic reference) or a footnote the source or sources from which you took the ideas.
\end{itemize}

A single standard of bibliographic references (citation) must be followed, from among the following alternatives:
\begin{itemize}
	\item APA (\url{http://pitt.libguides.com/c.php?g=12108\&p=64730})
	\item IEEE (\url{https://ieee-dataport.org/sites/default/files/analysis/27/IEEE\%20Citation\%20Guidelines.pdf})
	\item Harvard (\url{https://libweb.anglia.ac.uk/referencing/harvard.htm})
	\item With references numbered in alphabetical order or in the order of appearance in the text (for example, the numbered style used by some ACM publications - \url{https://www.acm.org/publications/authors/reference-formatting})
\end{itemize}

In LaTeX it is very easy to use references in a correct and consistent way, either by adding a
\verb!\begin{thebibliography}!
section (see at the end of this section), or through a separate bib file, using the command
\verb!\bibliography{}!,
as we do below by using the ``bibliography.bib'' file. In any case, in LaTeX you will need to use the command
\verb!\cite{}!
to add references, and this command must be used directly in the text, where you want the citation to appear, as in the following examples:
\begin{itemize}
	\item Journal article: ~\cite{article};
	\item Conference article:~\cite{proc};
	\item Book: ~\cite{book};
	\item Web link: ~\cite{silva};
\end{itemize}

\textbf{Important}: in this section, usually only the bibliographic entries appear (that is, just the listing of references). Citing them via the cite command and explanations related to them must be done in the previous sections. The citation above was made here only for illustration.

% This is how you specify the use of a file with bibliographic references:
\bibliographystyle{plain}
\bibliography{bibliography}

%% Another option would have been to include articles directly in this file
%% as follows:
%% \begin{thebibliography}{ABC}
%%
%% \bibitem{article}
%%  H. Baali, H. Djelouat, A. Amira and F. Bensaali,
%%  ``Empowering Technology Enabled Care Using IoT and Smart Devices:
%   A Review''. In: IEEE Sensors Journal, vol. 322 (10), pp. 891--921, 1905.
%%
%% (more \bibitem items here...)
%%
%% \end{thebibliography}

%% If you want a section to start on a new page, you can force this with the ``\newpage'' command, as below:

%\newpage

\chapter*{Appendices}\addcontentsline{toc}{chapter}{Appendices}

The appendices are optional.
What can go into the appendices:
\begin{itemize}
\item	An example of a configuration or build file;
\item	A table larger than half a page;
\item	A figure larger than half a page;
\item	A source code snippet larger than half a page;
\item	A set of screenshots;
\item	An example of running some commands plus their output;
\item 	The appendices contain things that take up more than one page and would interrupt the natural flow of reading the text.
\end{itemize}

\begin{appendices}

\chapter{Code Excerpts} % Introduces a new appendix
\ldots


\end{appendices}
\end{document}
